\addchap{Preface}

These notes began as an act of self-defence.

A Cambridge lecture moves at about the speed of the lecturer's handwriting, which is
much faster than the speed of understanding. You arrive, you copy down a definition, a
theorem and a proof, and an hour later you leave with a notebook full of statements
that are correct on the page and only approximately correct in your head. My response,
in my first term, was to go home in the evening and write the lecture again from
scratch: in my own words, in my own order, with the missing steps put back in wherever
I had lost the thread. It was much slower than reading. It was also the only thing that
reliably worked, because you cannot write a proof out honestly without discovering
exactly which line you did not follow.

I did this for one course, and then for another, and then it became simply what I did.
Three years later there were forty-four of these documents sitting in a folder,
covering every course I took for Parts IA, IB and II of the Mathematical Tripos,
somewhere close to three thousand pages in all. This book is that folder, bound
together and given a spine.

\section*{What this book is}

It is a single volume containing the complete set of my Cambridge undergraduate lecture
notes: eight courses from Part IA, fifteen from Part IB, and twenty-one from Part II\@.
Each chapter is one course, more or less exactly as I wrote it at the time. The
assembling is new; the mathematics is not.

That is worth being precise about, because it gives the book a texture you would never
choose if you sat down to write a textbook. The Tripos is a collection of courses, not
a single narrative, and courses overlap. Cauchy's theorem is proved carefully in
Complex Analysis and then used briskly in Complex Methods. The Gauss--Bonnet theorem
appears once in IB Geometry and again, from a greater height, in II Differential
Geometry. The Fourier transform turns up in four separate chapters wearing four
separate sets of hypotheses, and the Baire category theorem is proved twice. Notation
drifts between chapters, and a symbol defined in one chapter means nothing in the next.

I have deliberately left all of this alone. Smoothing it out would have meant rewriting
forty-four courses into one, which is a different book and not one I have written. What
you have instead is an honest record of what the subject actually looks like from
inside an undergraduate degree -- overlapping, repetitive in places, and occasionally
inconsistent with itself. The repetitions are, I think, among the more instructive
things here. Seeing the same theorem proved twice at two levels of generality tells you
something about the theorem that seeing it once does not.

\section*{What this book is not}

It is not a substitute for a textbook, and it is certainly not a substitute for going
to the lectures. Where a textbook would motivate at length, cite the literature and
give three proofs, these notes give one proof and move on. They were written under time
pressure by somebody meeting the material for the first time, and they show it.

Nor is this a light tour. The book that inspired the shape of this front matter --
Evan Chen's \emph{An Infinitely Large Napkin} -- sets out to give you the feeling of a
field in forty pages, and cheerfully omits any proof that does not teach something.
This book does close to the opposite, not because that approach is wrong but because
these notes were written for a different purpose. They are working notes: rigorous,
complete against the schedules, and in places genuinely terse. There are stretches that
are close to unreadable unless you already half know what is going on, which is what
happens when you write for the version of yourself who sat in the lecture that morning.

Finally, these are not transcriptions. Each chapter says so at its start, and each one
means it. I reordered material, replaced proofs I did not like, and added examples that
were never lectured. All lectured content should be covered, but the route through it
is frequently my own, and any error you find in a proof is far more likely to be mine
than the lecturer's.

\section*{Who it is for}

Mostly, honestly, for me -- somewhere to look things up, and a record of what three
years of work actually amounted to. But there are two other readers I had in mind while
putting it together.

The first is a Cambridge student sitting one of these courses now, who wants a second
account of the lectures to set beside their own notes, and would rather have all of
them in one file than in forty-four. The second is anyone learning this material from
outside Cambridge who wants to see what a complete and rigorous undergraduate path
through pure and applied mathematics looks like when it is written down end to end. The
Tripos is unusually broad -- there are not many degrees in which the same student
proves the Riemann--Roch theorem and computes the flow of a viscous fluid down a
slope -- and that breadth is most of what makes a document like this worth assembling
at all.

\section*{Debts}

The largest debt is to the lecturers. Almost every good idea in this book was explained
to me first by somebody standing at a blackboard, and the choice of material, the order
of the arguments and very often the proofs themselves are theirs. What I have added is
the writing-out, and occasionally the impertinence of thinking I could arrange it
better.

I owe a second debt to the other note-takers. Cambridge has a generous and slightly
anarchic tradition of students typing up their courses and putting them online for
everybody else, and I leaned on it constantly -- Dexter Chua's notes in particular were
a fixture of my degree, and there were plenty of others whose PDFs I went to when my
own record of a lecture had gone hopelessly wrong. The style file all of this is set in
descends, in its turn, from Evan Chen's and Tony Zhang's. Very little here was invented
from nothing, and that is rather the point.

\section*{Errors}

There are certainly some. If you find one -- a broken proof, a missing hypothesis, a
figure that says the opposite of the text -- I would genuinely like to hear about it,
at \texttt{ak2316@cam.ac.uk}. The notes are kept in a public repository as well, so
corrections can go there directly.

\bigskip
\noindent\emph{Adam Kelly}
